\documentclass[10pt,a4paper]{amsart}
\usepackage{amsmath,amssymb,amsthm,graphicx}
\usepackage{verbatim}
\newenvironment{reference}{\comment}{\endcomment}
%\newenvironment{reference}{}{}
\newenvironment{slogan}{\comment}{\endcomment}
\newenvironment{history}{\comment}{\endcomment}

% For commutative diagrams we use Xy-pic
\usepackage[all]{xy}

% We use 2cell for 2-commutative diagrams.
\xyoption{2cell}
% Theorem environments.
%
\theoremstyle{plain}
\newtheorem{theorem}[subsection]{Theorem}
\newtheorem{proposition}[subsection]{Proposition}
\newtheorem{lemma}[subsection]{Lemma}

\theoremstyle{definition}
\newtheorem{definition}[subsection]{Definition}
\newtheorem{example}[subsection]{Example}
\newtheorem{exercise}[subsection]{Exercise}
\newtheorem{situation}[subsection]{Situation}

\theoremstyle{remark}
\newtheorem{remark}[subsection]{Remark}
\newtheorem{remarks}[subsection]{Remarks}

\numberwithin{equation}{subsection}

% Macros
%
\def\lim{\mathop{\mathrm{lim}}\nolimits}
\def\colim{\mathop{\mathrm{colim}}\nolimits}
\def\Spec{\mathop{\mathrm{Spec}}}
\def\Hom{\mathop{\mathrm{Hom}}\nolimits}
\def\Ext{\mathop{\mathrm{Ext}}\nolimits}
\def\SheafHom{\mathop{\mathcal{H}\!\mathit{om}}\nolimits}
\def\SheafExt{\mathop{\mathcal{E}\!\mathit{xt}}\nolimits}
\def\Sch{\mathit{Sch}}
\def\Mor{\mathop{\mathrm{Mor}}\nolimits}
\def\Ob{\mathop{\mathrm{Ob}}\nolimits}
\def\Sh{\mathop{\mathit{Sh}}\nolimits}
\def\NL{\mathop{N\!L}\nolimits}
\def\CH{\mathop{\mathrm{CH}}\nolimits}
\def\proetale{{pro\text{-}\acute{e}tale}}
\def\etale{{\acute{e}tale}}
\def\QCoh{\mathit{QCoh}}
\def\Ker{\mathop{\mathrm{Ker}}}
\def\Im{\mathop{\mathrm{Im}}}
\def\Coker{\mathop{\mathrm{Coker}}}
\def\Coim{\mathop{\mathrm{Coim}}}

% Boxtimes
%
\DeclareMathSymbol{\boxtimes}{\mathbin}{AMSa}{"02}

%
% Macros for moduli stacks/spaces
%
\def\QCohstack{\mathcal{QC}\!\mathit{oh}}
\def\Cohstack{\mathcal{C}\!\mathit{oh}}
\def\Spacesstack{\mathcal{S}\!\mathit{paces}}
\def\Quotfunctor{\mathrm{Quot}}
\def\Hilbfunctor{\mathrm{Hilb}}
\def\Curvesstack{\mathcal{C}\!\mathit{urves}}
\def\Polarizedstack{\mathcal{P}\!\mathit{olarized}}
\def\Complexesstack{\mathcal{C}\!\mathit{omplexes}}
% \Pic is the operator that assigns to X its picard group, usage \Pic(X)
% \Picardstack_{X/B} denotes the Picard stack of X over B
% \Picardfunctor_{X/B} denotes the Picard functor of X over B
\def\Pic{\mathop{\mathrm{Pic}}\nolimits}
\def\Picardstack{\mathcal{P}\!\mathit{ic}}
\def\Picardfunctor{\mathrm{Pic}}
\def\Deformationcategory{\mathcal{D}\!\mathit{ef}}
\usepackage{cleveref}
\crefname{originalreference}{original reference}{original references}
\Crefname{originalreference}{Original reference}{Original references}
\makeatletter
\newlabel{lemma-nodal-vector-fields}{{0E68}{source}{}{sourcepointer-1}{}}
\newlabel{lemma-nodal-vector-fields@cref}{{[originalreference][0][]0E68}{[1][0][]source}{}{}{}}
\makeatother
\crefname{originalreference}{original reference}{original references}
\Crefname{originalreference}{Original reference}{Original references}
\makeatletter
\newlabel{lemma-smooth-vector-fields}{{0E67}{source}{}{sourcepointer-2}{}}
\newlabel{lemma-smooth-vector-fields@cref}{{[originalreference][0][]0E67}{[1][0][]source}{}{}{}}
\makeatother
\title{A nodal curve of genus at least 2 without rational tails or bridges has no global vector fields}
\author{The Stacks Project Authors}
\date{stacks-project:curves.tex:local-snapshot-20261012; GNU FDL 1.2 or later}
\begin{document}
\maketitle
\noindent\textit{Editorial note.} All delivered mathematical text is the author's own, copied from the pinned Stacks Project chapter file curves.tex; no statement or proof step was rewritten, reordered or removed.\par\medskip
\noindent\textit{Editorial note.} The counted proof is the author's complete argument: the normalization transport, the three component cases with their fixed-point counts, the exclusion of the glued, bridge, and tail configurations, and the concluding three-point vanishing are all retained verbatim.\par\medskip
\noindent\textit{Editorial note.} The label lemma-stable-vector-fields names the counted statement, and the proof environment immediately following it in the source is delivered verbatim as its proof.\par\medskip
\noindent\textit{Editorial note.} The counted statement is carried by the author's own lemma environment in the chapter source.\par\medskip
\noindent\textit{Allowed standard prerequisites (editorial).}\begin{itemize}
\item Derivations across the normalization at nodes. Same chapter, Lemma lemma-nodal-vector-fields. Transports a derivation D on X to a derivation D' on the normalization X\textasciicircum{}nu whose behaviour at the preimages of nodes mirrors D.
\item Vanishing of derivations on smooth projective curves. Same chapter, Lemma lemma-smooth-vector-fields. Part (1) kills derivations on components of genus > 1, part (2) on genus-1 components fixing a point, part (3) on rational components fixing three points.
\item The excluded configurations. Same chapter, Examples example-rational-tail and example-rational-bridge. Rule out the low-node configurations that would let a rational or elliptic component survive with a nonzero derivation.
\end{itemize}\medskip
\ifdefined\hypertarget\hypertarget{sourcepointer-1}{}\fi
\noindent\textit{Source reference.} Original reference 0E68: sources/stacks-curves/pinned-original-source.tex, source lines 6916--6916. Prerequisite pointer invoked at proof line(s) 6972; it names a standard fact of the same chapter that the counted argument uses, and every step of the counted proof built on it is retained verbatim in the delivered body.\par\medskip
\ifdefined\hypertarget\hypertarget{sourcepointer-2}{}\fi
\noindent\textit{Source reference.} Original reference 0E67: sources/stacks-curves/pinned-original-source.tex, source lines 6878--6878. Prerequisite pointer invoked at proof line(s) 6975; it names a standard fact of the same chapter that the counted argument uses, and every step of the counted proof built on it is retained verbatim in the delivered body.\par\medskip
\begin{lemma}
\label{lemma-stable-vector-fields}
Let $k$ be an algebraically closed field.
Let $X$ be an at-worst-nodal, proper, connected
$1$-dimensional scheme over $k$. Assume the genus of $X$ is at least
$2$ and that $X$ has no rational tails
or bridges. Then
$\text{Der}_k(\mathcal{O}_X, \mathcal{O}_X) = 0$.
\end{lemma}

\begin{proof}
Let $D \in \text{Der}_k(\mathcal{O}_X, \mathcal{O}_X)$.
Let $X^\nu$ be the normalization of $X$.
Let $D' \in \text{Der}_k(\mathcal{O}_{X^\nu}, \mathcal{O}_{X^\nu})$
be the element corresponding to $D$ via Lemma \ref{lemma-nodal-vector-fields}.
Let $C \subset X^\nu$ be an irreducible component.
If the genus of $C$ is $> 1$, then $D'|_{\mathcal{O}_C} = 0$
by Lemma \ref{lemma-smooth-vector-fields} part (1).
If the genus of $C$ is $1$, then there is at least one closed point $c$ of $C$
which maps to a node on $X$ (since otherwise $X \cong C$ would have genus $1$).
By the correspondence this means that
$D'|_{\mathcal{O}_C}$ fixes $c$ hence is zero by
Lemma \ref{lemma-smooth-vector-fields} part (2).
Finally, if the genus of $C$ is zero, then there are at least
$3$ pairwise distinct closed points $c_1, c_2, c_3 \in C$
mapping to nodes in $X$, since otherwise either
$X$ is $C$ with two points glued (two points of $C$ mapping
to the same node), or
$C$ is a rational bridge (two points mapping to different nodes of $X$), or
$C$ is a rational tail (one point mapping to a node of $X$).
These three possibilities are not permitted since
$C$ has genus $\geq 2$ and has no rational bridges, or rational tails.
Whence $D'|_{\mathcal{O}_C}$ fixes $c_1, c_2, c_3$ hence is zero by
Lemma \ref{lemma-smooth-vector-fields} part (3).
\end{proof}

\end{document}
